% Originally: new section (no predecessor in the week-based skeleton).
% Quelle: Linus Lüchinger/Slides/Slides-05-28.pdf, slides 2--3.

\section{Compactness in spaces of functions: Ascoli--Arzelà}
\label{sec:ascoli_arzela}

% Supplement: Linus Lüchinger/Slides/Slides-05-28.pdf, slides 2--3
\begin{ainote}
Corsin's notes stop at \cref{thm:heine_borel}, which settles compactness in $\mathbb{R}^n$ and
says nothing about spaces of \emph{functions}. Linus devotes his final session to the missing
theorem (Theorem 15.47 of the lecture notes), presenting it as \qt{the Bolzano--Weierstrass
theorem, one level up}. It is worth having here for a concrete reason: it is the compactness
input that \cref{rem:peano_vs_picard} appeals to without naming, when it says Peano's theorem
extracts a convergent subsequence from a family of approximate solutions.
\end{ainote}

\subsection{The space of continuous functions on a compact set}

\begin{definition}[Supremum norm on $C(K,\mathbb{R}^m)$]
\label{def:sup_norm_function_space}
Let $K \subseteq \mathbb{R}^n$ be compact. Write $C(K,\mathbb{R}^m)$ for the set of continuous
maps $f : K \to \mathbb{R}^m$, a real vector space under pointwise operations, and equip it with
the \newterm{supremum norm} \germanfor{Supremumsnorm}
\[\lVert f\rVert_\infty := \sup_{x \in K} \lvert f(x)\rvert.\]
\end{definition}

The supremum is finite because $\lvert f\rvert$ is continuous on a compact set, hence bounded
(\cref{sec:why_compactness_matters}), which is the same point already made for the supremum metric
in \cref{ex:sup_metric_is_metric}. Convergence in this norm is exactly \emph{uniform} convergence.

\begin{proposition}[$C(K,\mathbb{R}^m)$ is complete]
\label{prop:function_space_complete}
For $K \subseteq \mathbb{R}^n$ compact, $\big(C(K,\mathbb{R}^m), \lVert\cdot\rVert_\infty\big)$
is a complete normed vector space.
\end{proposition}

\begin{proof}
Let $(f_k)_{k\in\mathbb{N}}$ be Cauchy for $\lVert\cdot\rVert_\infty$. For each fixed $x \in K$
we have $\lvert f_k(x)-f_\ell(x)\rvert \leq \lVert f_k-f_\ell\rVert_\infty$, so $(f_k(x))_k$ is
Cauchy in $\mathbb{R}^m$ and converges, by completeness of $\mathbb{R}^m$, to a limit we call
$f(x)$. This defines a function $f : K \to \mathbb{R}^m$.

The convergence is uniform. Given $\varepsilon>0$ choose $N$ with
$\lVert f_k-f_\ell\rVert_\infty < \varepsilon$ for $k,\ell \geq N$; fixing $k \geq N$ and letting
$\ell \to \infty$ in $\lvert f_k(x)-f_\ell(x)\rvert < \varepsilon$ gives
$\lvert f_k(x)-f(x)\rvert \leq \varepsilon$ for \emph{every} $x$, i.e.\
$\lVert f_k - f\rVert_\infty \leq \varepsilon$.

Finally $f$ is continuous, being a uniform limit of continuous functions, by the standard
$\varepsilon/3$ argument from \emph{Analysis I}. Hence $f \in C(K,\mathbb{R}^m)$ and
$f_k \to f$ in the norm.
\end{proof}

\subsection{Why Bolzano--Weierstrass fails one level up}

In $\mathbb{R}^n$, bounded sequences have convergent subsequences. The naive transfer of that
statement to $C(K,\mathbb{R}^m)$ is \textbf{false}, and it is worth seeing exactly how.

% Generator: Claude Opus 5 (High)
\begin{aiexample}[A bounded sequence of functions with no convergent subsequence]
\label{ex:bounded_no_convergent_subsequence}
Take $K := [0,1]$, $m := 1$ and $f_k(x) := x^k$. Every $f_k$ is continuous with
$\lVert f_k\rVert_\infty = 1$, so the sequence is bounded, and indeed lies on the unit sphere.

Yet no subsequence converges in $C([0,1])$. Pointwise,
\[f_k(x) \longrightarrow g(x) := \begin{cases} 0, & 0 \leq x < 1, \\ 1, & x = 1,\end{cases}\]
and the same holds for every subsequence. If some $f_{k_j}$ converged uniformly to a limit
$f \in C([0,1])$, then it would converge pointwise to $f$ as well, forcing $f = g$; but $g$ is
not continuous. So no subsequence can converge.

\begin{remark}[The diagnosis]
Boundedness was never the issue. What goes wrong is that the $f_k$ become arbitrarily
\emph{steep} near $x=1$: to keep $\lvert f_k(x)-f_k(y)\rvert$ below a fixed $\varepsilon$ there,
one needs $\lvert x-y\rvert$ smaller and smaller as $k$ grows. Each individual $f_k$ is
uniformly continuous (\cref{def:uniformly_continuous}), being continuous on a compact interval.
What fails is that no single $\delta$ works for all of them at once. Ruling that out is precisely
what the next definition does.
\end{remark}
\end{aiexample}

\subsection{Equicontinuity, and the theorem}

\begin{definition}[Equicontinuity]
\label{def:equicontinuous}
A family $(f_k)_{k\in\mathbb{N}}$ in $C(K,\mathbb{R}^m)$ is \newterm{equicontinuous}
\germanfor{gleichgradig stetig} if
\[\forall \varepsilon>0\ \ \exists \delta>0\ \ \forall k \in \mathbb{N}\ \ \forall x,y \in K :
  \quad \lvert x-y\rvert < \delta \implies \lvert f_k(x)-f_k(y)\rvert < \varepsilon.\]
\end{definition}

% Feedback: Gemini 3.1 Pro (2026-08-11) --- called didactically outstanding: the three-way
% quantifier comparison packs the chapter's most abstract notion into three bullets. Keep the
% three-item list and the closing sentence that names the gap.
\begin{remark}[Reading the quantifiers]
\label{rem:equicontinuity_quantifiers}
Compare the three conditions, which differ only in what $\delta$ is allowed to depend on:
\begin{itemize}
  \item \textbf{Continuity} of one $f$: $\delta$ may depend on $\varepsilon$ \emph{and on the
    point} $x$.
  \item \textbf{Uniform continuity} of one $f$: $\delta$ may depend on $\varepsilon$ only, so
    there is one $\delta$ for all points.
  \item \textbf{Equicontinuity} of a family: $\delta$ may depend on $\varepsilon$ only, so there
    is one $\delta$ for all points \emph{and all members of the family}.
\end{itemize}
So equicontinuity is uniform continuity made uniform in one further variable, the index. Each
$x^k$ of \cref{ex:bounded_no_convergent_subsequence} is uniformly continuous, and the family is
not equicontinuous; that is the whole gap.
\end{remark}

% Quelle: Linus Lüchinger/Slides/Slides-05-28.pdf, slide 3
\begin{theorem}[Ascoli--Arzelà]
\label{thm:ascoli_arzela}
Let $K \subseteq \mathbb{R}^n$ be compact and let $(f_k)_{k\in\mathbb{N}}$ be a sequence in
$C(K,\mathbb{R}^m)$ that is
\begin{enumerate}[label=\textbf{(\alph*)}]
  \item \textbf{bounded}: $\sup_k \lVert f_k\rVert_\infty < \infty$; and
  \item \textbf{equicontinuous} in the sense of \cref{def:equicontinuous}.
\end{enumerate}
Then $(f_k)$ has a subsequence converging uniformly to some $f \in C(K,\mathbb{R}^m)$.
\end{theorem}

% Linus states the theorem without proof, remarking that it is "mostly just used in proofs".
\begin{ainote}
None of the transcribed notes prove this theorem; Linus states it and remarks that it is
\qt{mostly just used in proofs}. The sketch below is included because its mechanism, a
diagonal argument, explains why \emph{both} hypotheses are needed, which the statement alone does
not. It is a sketch and not a proof: the details are omitted.
\end{ainote}

\begin{remark}[Proof sketch: the diagonal argument]
\label{rem:ascoli_diagonal_sketch}
Fix a countable dense subset $D = \{q_1,q_2,\dots\} \subseteq K$ (the points of $K$ with rational
coordinates will do).

\textbf{Step 1: converge on $D$, using boundedness.} The values $(f_k(q_1))_k$ form a bounded
sequence in $\mathbb{R}^m$, so by Bolzano--Weierstrass some subsequence converges at $q_1$. Pass
to a further subsequence converging also at $q_2$, then at $q_3$, and so on. Each refinement is
legitimate, but no single one of them handles all of $D$. Taking the \emph{diagonal} sequence,
whose $j$-th member is the $j$-th member of the $j$-th subsequence, produces one subsequence
converging at every point of $D$ simultaneously.

\textbf{Step 2: upgrade to uniform convergence, using equicontinuity.} Convergence on a dense
set says nothing on its own about the rest of $K$; a sequence can converge at every rational and
oscillate wildly in between. Equicontinuity forbids that: given $\varepsilon$, pick the $\delta$
it supplies, cover the compact $K$ by finitely many balls of radius $\delta$ centred at points
of $D$, and control any $x$ by the nearby dense point. This makes the diagonal subsequence
uniformly Cauchy, hence convergent by \cref{prop:function_space_complete}.
\end{remark}

\begin{remark}[Where each hypothesis does its work]
The two hypotheses are not interchangeable, and the sketch shows which does what:
\textbf{boundedness} is what lets Bolzano--Weierstrass run at each individual point, and
\textbf{equicontinuity} is what converts pointwise control on a dense set into uniform control
everywhere. Delete the second and \cref{ex:bounded_no_convergent_subsequence} is a
counterexample; delete the first and the constant sequences $f_k \equiv k$ are equicontinuous
(they are all constant, so any $\delta$ works) with no convergent subsequence at all.
\end{remark}

% Generator: Claude Opus 5 (High)
\begin{aiexercise}[A uniform derivative bound gives compactness]
\label{ex:uniform_lipschitz_family}
Let $(f_k)_{k\in\mathbb{N}}$ be a sequence in $C^1([0,1],\mathbb{R})$ with
\[\lvert f_k'(x)\rvert \leq 1 \quad \text{for all } x \in [0,1] \text{ and all } k,
  \qquad\text{and}\qquad f_k(0) = 0 \text{ for all } k.\]
\begin{enumerate}[label=\textbf{(\alph*)}]
  \item Show that the family is equicontinuous.
  \item Show that it is bounded in $\lVert\cdot\rVert_\infty$.
  \item Conclude that some subsequence converges uniformly, and explain why the argument would
    break down if the bound $\lvert f_k'\rvert \leq 1$ were replaced by \qt{each $f_k$ is
    differentiable}.
\end{enumerate}
\exsol{sol:uniform_lipschitz_family}
\end{aiexercise}

% Transition: Claude Opus 5 (High)
That closes the account of compactness. What follows returns to $\mathbb{R}^n$ and to
differentiation, where compactness reappears as a hypothesis rather than a subject.
